\begin{frame}{Overview of Our Methodology: Pipeline}
% --- Professional Publication-Grade Technical Palette ---
\definecolor{techDark}{RGB}{33, 37, 41}      % Charcoal black for text and lines
\definecolor{techGray}{RGB}{108, 117, 125}   % Neutral slate gray for labels
\definecolor{techLight}{RGB}{248, 249, 250}  % Off-white for clean box backgrounds

\vspace{0.4cm}

\begin{center}
\begin{tikzpicture}[
    scale=0.9,
    every node/.append style={scale=0.9},
    node distance=0.72cm,
    % Main phase box style - matches your sketch dimensions precisely
    box/.style={
        rectangle, 
        draw=techDark!80, 
        thick, 
        fill=techLight, 
        text width=2.8cm, 
        align=center, 
        rounded corners=3pt, 
        inner sep=1.2ex, 
        font=\scriptsize, 
        minimum height=1.7cm
    },
    % Sub-boxes inside Step 2 (simulation and proof)
    subbox/.style={
        rectangle, 
        draw=techDark!80, 
        thick, 
        fill=white, 
        text width=1.3cm, 
        align=center, 
        rounded corners=2pt, 
        inner sep=0.8ex, 
        font=\scriptsize, 
        text=techDark, 
        minimum height=0.6cm
    },
    % Crisp, clean arrows matching technical drafts
    arrow/.style={
        ->, 
        >=Stealth, 
        thick, 
        draw=techDark!90
    }
]

    % --- Simple Technical Actor Icon (faithful to sketch) ---
    \tikzset{
      actor/.pic={
        \draw[thick, techDark] (0,0.18) circle (0.07); % head
        \draw[thick, techDark] (0,0.11) -- (0,-0.12); % body
        \draw[thick, techDark] (-0.11,0.02) -- (0.11,0.02); % arms
        \draw[thick, techDark] (0,-0.12) -- (-0.08,-0.3); % left leg
        \draw[thick, techDark] (0,-0.12) -- (0.08,-0.3); % right leg
      }
    }

    % ==================== PHASE 1 ====================
    \node[box] (step1) {\textbf{1. Generate candidate algorithms}};

    % Input: Scenario Actor (stands above Stage 1)
    \coordinate (user1) at ([yshift=0.85cm]step1.north);
    \pic at (user1) {actor};
    \node[above=0.25cm of user1, font=\scriptsize, text=techDark] (scenario) {Goals};
    \draw[arrow] ([yshift=-0.3cm]user1) -- (step1.north);


    % ==================== PHASE 2 ====================
    % Outer validation container
    \node[box, right=of step1, minimum width=4.0cm] (step2) {};
    
    % Step 2 Internal Header
    \node[anchor=north, font=\scriptsize\bfseries, text=techDark, yshift=-0.12cm] 
        at (step2.north) {2. Validation};
    
    % Nested simulation & proof boxes
    \node[subbox, anchor=south, xshift=-0.9cm, yshift=0.18cm] (sim) at (step2.south) {\strut simulation};
    \node[subbox, anchor=south, xshift=0.9cm, yshift=0.18cm] (proof) at (step2.south) {\strut proof};
    
    % Conjunction symbol (+) between simulation and proof
    \node[font=\scriptsize\bfseries, text=techDark!70, anchor=center, yshift=0.48cm] 
        at (step2.south) {+};

    % Actor pointing directly up to the proof sub-box
    \coordinate (user2) at ([yshift=-0.9cm]proof.south);
    \pic at (user2) {actor};
    \draw[arrow] ([yshift=0.25cm]user2) -- (proof.south);


    % ==================== PHASE 3 ====================
    \node[box, right=of step2] (step3) {\textbf{3. Classification}};


    % ==================== INTER-PHASE DATA FLOWS ====================
    % C-labeled transit arrow
    \draw[arrow] (step1.east) -- node[above=0.02cm, font=\small] {$\mathcal{C}$} (step2.west);

    % A-labeled transit arrow
    \draw[arrow] (step2.east) -- node[above=0.02cm, font=\small] {$\mathcal{A}$} (step3.west);

    % Output arrow (down-then-right layout matching sketch output style)
    \draw[arrow] (step3.east) -- ++(0.3, 0) -- ++(0, -1.0) 
        node[below=0.05cm, font=\scriptsize, text=techDark, align=center] {$\mathcal{A}$ \\ + classes};

\end{tikzpicture}
\end{center}

\end{frame}
